\section{Companion lab: evaluate an expression language}
\label{revision-lab:expression}

CBOR makes nested structure explicit in its heads and container counts. A text
expression instead uses punctuation, operator rules and a byte cursor. This lab
asks you to implement a small language whose result depends on those rules, not
merely to replace characters with arithmetic instructions. It revises original
project \texttt{017-expression-evaluator}.

Use \texttt{expression-evaluator.asm} in \texttt{revision/labs/files} as your
working skeleton. Implement \texttt{eval\_expr}; the entry deliberately stops with
\texttt{break} until replaced. Keep the supplied caller and data arenas. The
separate complete \texttt{expression-arithmetic.asm} demonstrates checked
operations without supplying the parser. Original grammar and feedback remain
in \texttt{packages/017-expression-evaluator}.

The prerequisites are bounded byte access, signed-word representation, full-width
function frames and the parent-state reasoning in the nested-data chapter.
The new work connects grammar to control flow, retains pending operations across
calls, and reports the first failed requirement precisely. These activities use
ordinary Run and Debug; their native guided lesson remains unfinished.

\subsection*{Give the input a grammar before writing the loop}

For \texttt{18 - 3 * (2 + -4)}, first evaluate the group to minus two, then the
product to minus six, then subtract that product from eighteen. The result is
24. A loop that immediately applies each operator to the next number would
instead reach fifteen before seeing multiplication. Operator precedence tells
you which larger operand must complete before an operation can run.

\begin{figure}[htbp]
  \centering
  \begin{minipage}{.94\linewidth}
    \begin{verbatim}
expression := term { ('+' | '-') term }
term       := unary { ('*' | '/') unary }
unary      := ('+' | '-') unary | primary
primary    := digits | '(' expression ')'
digits     := one or more ASCII decimal digits
\end{verbatim}
  \end{minipage}
  \caption{The exercise language. Braces mean zero or more repetitions; the vertical bar separates alternatives. Quoted characters are actual input bytes. Parentheses outside quotes group the grammar notation. Whitespace may occur between tokens, but cannot join two digit sequences.}
  \label{fig:revision-expression-grammar}
\end{figure}


Read the expression rule as: parse one term, then repeatedly accept a plus or
minus followed by another complete term. The term rule does the same for
multiplication and division, using unary values as operands. A unary sign owns
a whole following unary value, so \texttt{--5} means the negation of minus five.
A primary is either a decimal literal or a group containing another expression.
The group owns both parentheses; the inner expression does not consume its
caller's closing parenthesis.

The loops implement left association. In \texttt{9-6/2-1}, the first right term
is \texttt{6/2}, giving three; the expression loop computes \(9-3=6\), then
\(6-1=5\). A recursive call for an unrestricted expression after each minus
could consume the rest as its right operand and change that association.
One routine per precedence level is a useful recursive-descent design, not a
required private organization.\footnote{For background on grammar-driven
recursive descent, see Robert Nystrom, \emph{Crafting Interpreters},
\href{https://craftinginterpreters.com/parsing-expressions.html}{Parsing expressions}.
This exercise defines its own grammar, numeric bounds and diagnostic order.}

Accept only ASCII digits, the six operator/grouping characters, and space,
tab, carriage return or line feed between tokens. The input is a selected byte
slice, not a zero-terminated string. An embedded zero is invalid. A vertical tab
is also invalid, even if another language calls it whitespace. \texttt{12 34}
contains two literals without an intervening operator; whitespace does not
turn it into 1234. No identifiers, comments, implicit multiplication, hexadecimal
or floating-point forms belong to this language.

\subsection*{Make the cursor and suspended work explicit}

Represent the input by a full base address, current byte offset and authoritative
length. Before any load, establish that the complete load lies in the slice.
An unsigned byte load is enough to distinguish ASCII punctuation and digits.
Reaching the end while skipping whitespace is ordinary progress; reaching it
while waiting for a primary or closing parenthesis is a syntax failure.

A useful private helper result is a value together with an advanced shared
cursor and an error state. When a term calls a unary parser, it must retain its
left value, pending operator and that operator's source offset. The recursive
child advances the cursor. Restoring the old cursor on return would undo the
child's progress and can repeat the same token indefinitely.

\begin{table}[htbp]
  \centering
  \begin{tabular}{r l l}
    \toprule
    Cursor offset & Completed work                & Pending operation             \\
    \midrule
    2             & Literal 18                    & Outer subtraction begins at 3 \\
    6             & Literal 3                     & Multiplication begins at 7    \\
    11            & Group's left term 2           & Addition begins at 12         \\
    16            & Unary minus has produced $-4$ & Group computes $2+(-4)$       \\
    17            & Group closes with value $-2$  & Term computes $3(-2)$         \\
    17            & Term returns $-6$             & Expression computes $18-(-6)$ \\
    \bottomrule
  \end{tabular}
  \caption{A trace of \texttt{18 - 3 * (2 + -4)}. Offsets are zero-based byte positions; 17 is the input end. A returned value and the cursor are different outputs of parsing. Whitespace skipping can advance the cursor between these boundaries.}
  \label{tab:revision-expression-trace}
\end{table}


The term awaiting the group in the worked input retains left value three,
operator \texttt{*} and offset seven. The group creates a nested expression
whose pending addition retains two and offset twelve. The unary parser remembers
the minus sign at fourteen while reading literal four. After that child
returns, negation gives minus four; the group gives minus two; the waiting
multiplication can finally run. Each frame owns its pending operation while
the cursor records shared forward progress.

Distinguish the public calling convention from private parser state. The public
\texttt{eval\_expr} must preserve every bit of \texttt{s0} through \texttt{s7},
\texttt{fp} and \texttt{sp}, and return through its incoming \texttt{ra}. A private
helper convention may deliberately share some saved registers as parser state,
provided the public wrapper restores the original caller values. Document that
convention; mechanically restoring every shared parse register in every private
helper would erase progress or the first error. Save return addresses and live
operands in full-width, balanced, properly aligned frames.

Another valid design uses a precedence threshold and one binary-expression
routine. After an operator of precedence $p$, parse its right operand with a
threshold excluding another operator of that same precedence; the enclosing
loop then performs equal-precedence operations left to right. This is often
called precedence climbing. Choose a design you can trace. Equivalent private
helpers and explicit stacks are accepted when they preserve the whole contract.

\subsection*{Compute widely, then prove the signed-word range}

Decimal literals are limited to zero through 2147483647. For current literal
value $n$ and next digit $d$, form $10n+d$ in a doubleword and test the result
before accepting that digit. The previous accepted value was at most 2147483647,
so this one-digit candidate fits comfortably in 64 bits. Record the current
digit's offset if it exceeds the positive limit. Long leading zeros are legal
within the 64-byte input bound.

Every intermediate arithmetic value must fit signed 32 bits, from minus
2147483648 through 2147483647. A wide candidate prevents word arithmetic from
hiding overflow before the test. Addition and subtraction of two word values
fit a signed doubleword; their product has magnitude at most \(2^{62}\).
After excluding zero division, a wide quotient also fits. The special quotient
\(-2147483648/-1=2147483648\) is a valid wide result that must fail the narrower
language range.

\Needspace{17\baselineskip}
Run the complete arithmetic example and break at \texttt{product\_ready}, then
\texttt{arithmetic\_ready}. Its first call computes \(3(-2)=-6\); the second
computes \(18-(-6)=24\). Both return zero status. At
\texttt{arithmetic\_stored}, the three native words hold product minus six,
status zero and result 24. The caller does not print them. Both supplied callers round the initial stack
pointer down to a sixteen-byte boundary before their first call.

\begin{figure}[!ht]
  \centering
  \begin{tikzpicture}[alt={Native Memory observation of the complete public checked-arithmetic example. An outline marks three 32-bit words: product minus six, zero status and result 24. These correspond to three times minus two, followed by eighteen minus that product.}]
    \node[anchor=north west,inner sep=0,font=\small\bfseries] (title) at (0,0) {A checked product feeds the next operation};
    \node[anchor=north west,inner sep=0] (a) at ([yshift=-2mm]title.south west) {\MARSRevisionCrop{revision/captures/expression/expression-arithmetic.png}{0.99\linewidth}{0.03125}{0.33888889}{0.58203125}{0.60888889}};
    \begin{scope}[shift={(a.south west)},x={($(a.south east)-(a.south west)$)},y={($(a.north west)-(a.south west)$)}]
      \draw[black,line width=.7pt] (.20,.005) rectangle (.77,.447);
    \end{scope}
  \end{tikzpicture}
  \caption{The first three little-endian words hold product minus six, status zero and final result 24. They are stored at \texttt{arithmetic\_stored} in the complete public arithmetic example. Native LaTeX supplies the title, crop and outline.}
  \label{fig:revision-expression-native-arithmetic}
\end{figure}


\Needspace{14\baselineskip}
In \texttt{checked\_binary}, \texttt{a0/a1} contain signed-word operands,
\texttt{a2} is the ASCII code for one of \texttt{+-*/}, and \texttt{a3} is that
operator's byte offset. It returns status in \texttt{v0} and value or error
offset in \texttt{v1}. The complete helper checks zero before executing
\texttt{ddiv}. Once \texttt{t0} holds a wide arithmetic candidate, its range
comparison begins:
\Needspace{8\baselineskip}
\MARSCodeLabel{Compare a wide candidate with its sign-extended low word}
\begin{minted}{mars-idea}
    sll $t1, $t0, 0
    beqc $t1, $t0, checked_value
    nop
\end{minted}

Here the word-sized \texttt{sll} intentionally discards the upper half and
sign-extends the low word. Equality with the original doubleword proves that
no information needed for a signed-word value was lost. It is a test using a
copy, not permission to replace the original candidate with a wrapped word.
For 2147483648, the copy is minus 2147483648, so equality fails.

Use \texttt{revision/labs/checks/expression-arithmetic.kts} to check the worked
helper and caller. Its operands and operators are already supplied as numbers;
it neither parses a token nor diagnoses malformed text. Your evaluator may
reuse or adapt it, with an explicit private interface, or implement equivalent
checked arithmetic differently.

\subsection*{Treat diagnostics as part of evaluation}

Return status zero and the signed value on success. Failure status minus one
means syntax, minus two overflow, and minus three division by zero. On failure,
\texttt{v1} changes meaning: it is a zero-based byte offset. The offset can equal
the input length when an item is missing at the end. Do not confuse a successful
numeric result such as ten with an error offset ten; inspect status first.

\begin{table}[htbp]
  \centering
  \begin{tabular}{l r r l}
    \toprule
    Selected input            & Status & Offset & First failed requirement       \\
    \midrule
    \texttt{2147483647+1?}    & $-2$   & 10     & Addition's word range          \\
    \texttt{2147483647+1*?}   & $-1$   & 13     & Multiplication's right operand \\
    \texttt{1/(2-2)}          & $-3$   & 1      & Completed denominator is zero  \\
    \texttt{1/(2-2}           & $-1$   & 6      & Group needs its closing byte   \\
    \texttt{-2147483648}      & $-2$   & 10     & Literal exceeds positive bound \\
    \texttt{-(-2147483647-1)} & $-2$   & 0      & Outer negation's word range    \\
    \bottomrule
  \end{tabular}
  \caption{Error identity follows the specified parse-and-evaluate procedure. A later invalid byte does not replace an earlier arithmetic failure; an operation cannot fail before its required operand has successfully completed.}
  \label{tab:revision-expression-errors}
\end{table}


In \texttt{2147483647+1?}, the right term one has completed, and the addition
fails at offset ten before final consumption checks the question mark. With
\texttt{2147483647+1*?}, that right term cannot finish: multiplication still
needs its right unary operand at thirteen. The syntax failure therefore occurs
before the outer addition can be attempted. A whole-input syntax prevalidation
pass would change the required diagnosis in the first example.

The following diagnostic figures use a completed evaluator with the same public
caller and data layout as your skeleton. At \texttt{expression\_stored}, the
row begins with length, status and value-or-offset, followed by a guard.

\begin{figure}[!ht]
\centering
\begin{tikzpicture}[alt={Actual native Memory after a completed evaluator processes 2147483647+1?. The row holds length thirteen, overflow status minus two, byte offset ten and an unchanged guard. An outline marks status and offset.}]
\node[anchor=north west,inner sep=0,font=\small\bfseries] (title) at (0,0) {Overflow precedes the trailing invalid byte};
\node[anchor=north west,inner sep=0] (a) at ([yshift=-2mm]title.south west) {\MARSRevisionCrop{revision/captures/expression/expression-overflow.png}{0.99\linewidth}{0.03125}{0.33888889}{0.58203125}{0.60888889}};
\begin{scope}[shift={(a.south west)},x={($(a.south east)-(a.south west)$)},y={($(a.north west)-(a.south west)$)}]
\draw[black,line width=.7pt] (.39,.005) rectangle (.77,.447);
\end{scope}
\end{tikzpicture}
\caption{The caller stores selected length thirteen, status minus two and operator offset ten. The outlined status/offset pair identifies the addition overflow in 2147483647+1?, before complete-consumption checking reaches the question mark. Native LaTeX supplies the title, crop and outline.}
\label{fig:revision-expression-native-overflow}
\end{figure}


Likewise, \texttt{1/(2-2} does not provide a denominator. The inner subtraction
has produced zero, but the group still needs its closing parenthesis at offset
six. Division cannot inspect that provisional group value as a completed operand.
After the closing parenthesis is present, the same denominator completes and
zero division is reported at the slash, offset one.

\begin{figure}[!ht]
  \centering
  \begin{tikzpicture}[alt={Actual native Memory after a completed evaluator processes the six bytes 1/(2-2. The row holds length six, syntax status minus one, byte offset six and an unchanged guard. An outline marks status and offset, distinguishing this from zero division.}]
    \node[anchor=north west,inner sep=0,font=\small\bfseries] (title) at (0,0) {An incomplete group is not a denominator};
    \node[anchor=north west,inner sep=0] (a) at ([yshift=-2mm]title.south west) {\MARSRevisionCrop{revision/captures/expression/expression-missing-close.png}{0.99\linewidth}{0.03125}{0.33888889}{0.58203125}{0.60888889}};
    \begin{scope}[shift={(a.south west)},x={($(a.south east)-(a.south west)$)},y={($(a.north west)-(a.south west)$)}]
      \draw[black,line width=.7pt] (.39,.005) rectangle (.77,.447);
    \end{scope}
  \end{tikzpicture}
  \caption{For the six-byte input 1/(2-2, the stored status is minus one and the offset is six, the missing closing parenthesis at the input end. The division has no completed denominator yet. Native LaTeX supplies the title, crop and outline.}
  \label{fig:revision-expression-native-missing-close}
\end{figure}


The leading minus in \texttt{-2147483648} is a unary operator, not part of the
literal. Its operand literal exceeds 2147483647 at the final digit, offset ten,
so the evaluator never reaches negation. Use \texttt{(-2147483647-1)} to obtain
the minimum through valid literals and an in-range subtraction. Negating that
completed value fails at the outer sign. An eventual multiplication by zero
does not excuse an earlier overflow in \texttt{(2147483647+1)*0}.

Keep the first error while returning through parents. A child failure prevents
its parent's pending arithmetic or closing-byte check from replacing that error.
Only a successful outer expression proceeds to final whitespace skipping and
complete-consumption checking. Useful offset records belong to different events:
a literal remembers the current digit; unary remembers its sign; binary remembers
its operator; missing syntax uses the current cursor or the input end.

\subsection*{Build and investigate your complete evaluator}

The public entry receives a full input pointer in \texttt{a0} and a length from
zero through 64 in \texttt{a1}. The slice lies in a 96-byte arena and may be
unaligned. All caller data, including input and guards, are read-only even for
temporary stores. Read only the selected slice among caller objects. Your own
constant tables and private stack may be used. Produce no console output.
The preserved caller source contains the worked 17-byte expression; it stores
the returned status and result after \texttt{expression\_ready}.

Implement and test boundaries that you can explain. A whitespace skip must
terminate at the slice end. A literal must stop before a nondigit and identify
the first excessive digit. A group must retain its value while checking its
closing byte. A unary sign must wait for a complete operand. A term must finish
multiplication and division before an expression loop uses its value. Connect
those pieces under the public saved-state wrapper, then require complete input
consumption. Intermediate helper tests support construction; they are not final
acceptance of the language.

After your evaluator is complete, break at \texttt{expression\_ready} for the
worked input and expect \texttt{v0}=0, \texttt{v1}=24. At
\texttt{expression\_stored}, inspect its stored words before exit overwrites the
service selector. The native figures use those shared caller boundaries in a
completed implementation; your private helpers need not match its layout.
Change only the selected input and length for each diagnostic case, using a
fresh copy or resetting the machine between ordinary Debug runs.

\Needspace{7\baselineskip}
Run the original public kernel and caller checks against your working source:
\begin{verbatim}
java -jar /path/to/cli.jar grade \
  packages/017-expression-evaluator/checks/kernel.kts \
  /path/to/your/expression-evaluator.asm
\end{verbatim}
Repeat with \texttt{program.kts}. The unfinished skeleton is expected to fail.
The kernel accepts different parsing strategies; it does not prescribe helper
names, frame sizes, traversal order or instruction counts. The caller checks
add observation and stored-result behavior. Passing finite cases does not prove
all 64-byte inputs.

Retain a prediction and explanation for \texttt{9-6/2-1}, consecutive unary
signs, separated digit sequences, the two question-mark examples and the two
denominator groups. Then call the function again after an error using another
input arena. A successful previous call must not seed the next parse, and a
previous failure must not remain as its error state. Compare return status,
value or byte offset, source preservation and balanced caller state; a program
that obtains the right arithmetic value while reading a guard is still incorrect.
